\documentclass[11pt]{article}
\usepackage{amsmath,amssymb,amsthm,hyperref}
\newtheorem{proposition}{Proposition}
\newcommand{\E}{\mathbb E}
\newcommand{\Pp}{\mathcal P}
\title{Actual-corner minors, nuclear tails, and the limit of matrix-only gluing}
\author{Campaign 3395: Engineer continuation}
\date{November 4, 2026}
\begin{document}
\maketitle

\paragraph{Status and comparison.}
The averaged circuit-specific nuclear-tail charge and the unrestricted
fixed-factor lower bound remain unproved. This note proves the actual-corner
minor budget, the requested truncated-SVD transfer (also in the reverse
direction), and a precise obstruction to a \emph{matrix-only} conversion.
The obstruction does not refute the charge under the full original sampler
with its polynomial-in-original-size exponential remainder.
Reviewed V14 remains the strongest achieved approximation theorem, at factor
$1+2^{-(2q+\lceil k_0\rceil+14)}$, not a fixed factor.
No verified novelty or publishable significance is claimed here.

Reuse reviewed K1--K2, J5/J8--J12, T1--T4, V8/V10/V14,
CD1--CD20, Z6--Z20, A4/A7, the compulsory cut and the original sampler.
Existing independent source attribution is unchanged. All new matrix
calculations below are proved directly; in particular, the incidence
matrix is a classical elementary construction, not a claimed new
finite-field or spectral theorem. No external gluing theorem is assumed.

\section{A minor bound for the actual four corners}
Retain CD5's probability law $\mu$, including the original pre-bit law,
the original uniform bits, and independent flips of probability $1/(2q)$.
For actual cut $Q_j=A_jV_j$, its honored row and column flips give
four actual one-hot corners with union $U_j\le U^*$.
Let $Z_j$ be the two-square matrix of $Q_j$ at those corners. Define
\[
 \phi(Z)=
 \begin{cases}
  \det(Z)^2/(\sum Z)^3,&\sum Z>0,\\
  0,&Z=0.
 \end{cases}                                                   \tag{RG1}
\]
This is not a minor of a globally completed error matrix.

\begin{proposition}
For any nonnegative two-square matrix $Z$ and any rank-at-most-one
matrix $R\ge Z$ entrywise,
\[
                    \phi(Z)\le\sum(R-Z).                       \tag{RG2}
\]
Consequently, for the actual local K1 completions,
\[
 \phi(Z_j)\le D_{Q_j}(U_j),\qquad
 \sum_j\E_\mu\phi(Z_j)\le e^3 b_Q.                              \tag{RG3}
\]
All coefficients and actual cut terms are retained.
\end{proposition}
\begin{proof}
Put $t=\sum Z$, $d=\sum(R-Z)$. The case $t=0$ is immediate.
Let $\sigma_1\ge\sigma_2\ge0$ be the singular values of $Z$.
The Frobenius norm is at most $t$, so $\sigma_1,\sigma_2\le t$.
Choose a unit vector in the kernel of $R$. Evaluating $Z$ on
that vector gives
$\sigma_2\le\|Z-R\|_{\mathrm F}\le d$.
Therefore
\[
 \phi(Z)=\frac{\sigma_1^2\sigma_2^2}{t^3}
       \le\frac{\sigma_2^2}{t}\le\sigma_2\le d.
\]
This argument includes $R=0$ and singular $Z$ and makes no division
by an entry, determinant, or completion error.

Use the K1 completion at $U_j$, restricted to the four corners.
It is a rank-one upper completion across the same compulsory partition.
If the two row strings or the two column strings coincide, $Z_j$ has
repeated rows or columns and its determinant is zero. No fourfold count
of a repeated corner is then used. Otherwise the injective original
bit-to-assignment map makes all four corners distinct. K2 bounds the
sum of their nonnegative errors by $D_{Q_j}(U_j)$, proving the first
inequality in (RG3). Finally sum over actual cuts under precisely
CD5's law and use CD7:
\[
 \sum_j\E_\mu\phi(Z_j)
 \le\sum_j\E_\mu D_{Q_j}(U_j)
 \le\E_\mu D_Q(U^*)\le e^3b_Q.
\]
There is one joint whole-output defect budget, not a sum of gate maxima.
\end{proof}

\section{The singular-value tail is an equivalent remaining obligation}
For every original conditional matrix, write
\[
 M_j=M_j(Q_j),\quad L_j=M_j(H_j),\quad W_j=M_j(W_j),\quad
 M_j=L_j+W_j,\quad a_j=m_j^{-2}\sum W_j,\quad m_j=2^{k_j}.
\]
In particular $a_j$ is still the actual \emph{heavy} conditional mean.
Set
\[
 I_j=1_{\{k_j\ge5k_0/4,\ a_j\ge2^{-k_0}\}},\qquad
 r_j=\lfloor2^{k_j/8}\rfloor,\quad
 \tau_j=\sum_{i>r_j}\sigma_i(M_j),\qquad
 {\cal T}(Q)=\sum_j\E\frac{I_j\tau_j^2}{m_j^3a_j}.              \tag{RG4}
\]
Every quotient is defined to be zero outside its displayed indicator;
there is no division by a zero heavy mean. The expectation is the
unchanged original outer law and each cut's original unhonored-bit fixing.
The truncations below are analytic matrices, not extracted circuits.

\begin{proposition}
For every allowed unrestricted circuit and admissible $N\ge2^{142}$,
\[
 {\cal T}(Q)\le{\cal E}_{\rm thick}(Q)
 \le3{\cal T}(Q)+54s\,2^{-3k_0/32}
                         +6s\,e^{-5k_0/32}.                    \tag{RG5}
\]
\end{proposition}
\begin{proof}
J5 gives $\operatorname{rank}L_j\le D_h$. If $h>0$, J8 gives
$\log D_h\le153k_0/2^{16}$; if $h=0$, $D_h=1$.
On $I_j=1$,
\[
 \frac{k_j\log2}{8}\ge\frac{5k_0}{64}
                         >\frac{153k_0}{2^{16}},
\]
using $\log2\ge1/2$. Since $D_h$ is an integer,
$D_h\le\lfloor2^{k_j/8}\rfloor=r_j$.

For completeness, the nuclear-norm best-rank-$r$ error of any matrix
$M$ is $\sum_{i>r}\sigma_i(M)$. Indeed, if $\operatorname{rank}B\le r$,
restrict the domain to $\ker B$, of codimension at most $r$.
The $i$th singular value of that restriction of $M$ is at least
$\sigma_{i+r}(M)$: intersect the span of its first $i+r$ right singular
vectors with $\ker B$ to get an $i$-dimensional test subspace.
Right multiplication by an orthogonal projection contracts nuclear
norm. Hence $\|M-B\|_*\ge\sum_{i>r}\sigma_i(M)$.
The truncated SVD attains equality. Applying this with $B=L_j$
gives $\tau_j\le\|W_j\|_*$, proving the first part of (RG5).

Let $M_j^{(r)}$ be its truncated SVD. The exact identity
$W_j=(M_j-M_j^{(r)})+M_j^{(r)}-L_j$ gives
\[
 \|W_j\|_*^2
 \le3\tau_j^2+3\|M_j^{(r)}\|_*^2+3\|L_j\|_*^2
 \le3\tau_j^2+27(r_j+D_h)m_j^2.                         \tag{RG6}
\]
Here $0\le L_j\le M_j\le3$ entrywise, so both relevant Frobenius
norms are at most $3m_j$; SVD truncation does not increase that norm.
On $I_j=1$, the normalized remainder is at most
\[
 \frac{27(r_j+D_h)}{m_ja_j}
 \le54\,2^{k_0-7k_j/8}\le54\,2^{-3k_0/32}.              \tag{RG7}
\]
Average, sum at most $s$ actual cuts, and add exactly V8's strip
$k_0\le k_j<5k_0/4$, whose contribution is at most
$6s e^{-5k_0/32}$. This proves (RG5).
\end{proof}

Thus a bound
\[
 {\cal T}(Q)\le Cb_Q+K M_n^\alpha(s+1)^\beta e^{-cN}             \tag{RG8}
\]
would give T5 with constants
$3C$, $3K+60$, $\alpha$, $\max(\beta,1)$, and
$\min(c,3/2^{140})$. The last rate uses $\log2\ge1/2$ in
(RG5). The reviewed fixed factor would be
$1+1/(768(3C+1))$. Conversely T5 implies (RG8) with its same
constants. These implications do not establish (RG8).

The known CD18--CD20 and Z6--Z20 witnesses do not obstruct (RG8):
their \emph{full actual cut-value matrices} have rank at most one
on every history, so $\tau_j=0$. This uses their displayed
separation across the compulsory partition, not a new replay.
The positive exponential terms in (RG5) are retained; they are not
silently set to zero.

\section{A density obstruction to matrix-only gluing}
The following all-size family rules out a conversion using only
the entry bound, low-band rank, heavy mean, and the averaged local
minor energy. It does not impose the full original outer sampler
or a small original-DAG realization.

Let $t=4^\ell$, $\ell\ge1$, and take a field $\mathbb F_t$.
Its existence needs no additional conjecture: in a splitting field
over $\mathbb F_2$, the $t$ distinct roots of $X^t-X$ form a field.
They are closed under addition, multiplication and nonzero inverses
by the characteristic-two power identity, and the derivative is $-1$.
A splitting field is obtained by successively adjoining roots using
quotients by irreducible factors.
Index both rows $(x,y)$ and columns $(a,b)$ by $\mathbb F_t^2$ and set
\[
 B_{(x,y),(a,b)}=1_{\{y=ax+b\}},\quad
 M=\tfrac12J+B,\quad L=\tfrac12J,\quad W=B,\quad
 m=t^2=2^k,\quad k=4\ell,\quad a_{\rm heavy}=1/t.        \tag{RG9}
\]
Each row and column of $B$ has exactly $t$ ones.
We have $0\le L\le M\le3/2$ and $\operatorname{rank}L=1$.
With $k_0=3\ell$, both $k\ge5k_0/4$ and
$a_{\rm heavy}\ge2^{-k_0}$ hold. Arbitrarily large such parameters
meet $k_0\ge256$.

\begin{proposition}
For $r=\lfloor m^{1/8}\rfloor=\lfloor t^{1/4}\rfloor$,
\[
 \frac{\tau_r(M)^2}{m^3a_{\rm heavy}}
       =\left(1-\frac1t-\frac{r-1}{t^2}\right)^2
       \longrightarrow1.                                    \tag{RG10}
\]
For \emph{any} four-corner coupling whose individual corners are
uniform on the row-column index set, including every sparse bit-flip
coupling used here,
\[
                       \E\phi(Z)\le\frac1{8t}.                \tag{RG11}
\]
Consequently no fixed multiple of this minor energy plus a
polynomial-in-$k$ exponentially decaying remainder controls (RG10).
\end{proposition}
\begin{proof}
Two distinct points with the same $x$ coordinate have no common
incident line; points with different $x$ coordinates have exactly one.
Thus, grouping rows by $x$,
\[
 BB^{\mathsf T}=tI+(J_t-I_t)\otimes J_t.
\]
Its eigenspaces are: constants, eigenvalue $t^2$; functions constant
inside each $x$ block and summing to zero across blocks, eigenvalue
zero and dimension $t-1$; and functions summing to zero inside
each block, eigenvalue $t$ and dimension $t(t-1)$.
The constant term in $M$ changes only the first singular value.
Its complete singular spectrum is therefore
\[
             (m/2+t)^1,\quad(\sqrt t)^{\,t(t-1)},\quad0^{\,t-1}.
                                                               \tag{RG12}
\]
Removing the first $r$ values gives (RG10).

For the four entries of $B$ let $w$ be their number of ones,
counted with their positions even if corners repeat.
All possibilities for $\phi(\tfrac12J+B_{\rm local})$ are
\[
\begin{array}{c|ccccc}
w&0&1&2\ {\rm adjacent}&2\ {\rm opposite}&3\quad 4\\ \hline
\phi&0&1/108&0&1/16&9/500\quad 0 .
\end{array}
\]
In each case $\phi\le w/32$. Each corner has incidence probability
$1/t$, so $\E w=4/t$, without an independence assumption between
corners. This proves (RG11). For any fixed $C,K,c>0$ and $d\ge0$,
$C/(8t)+Kk^d e^{-ck_0}\to0$, whereas (RG10) tends to one.
The same conclusion holds when $q$ is any fixed constant multiple
of $k$; in particular the prescribed $q/k_0$ scale does not rescue
this matrix-only estimate.
\end{proof}

This is not the previous rank-one completion-error obstruction.
Here $W$ has rank $t^2-t+1$, and truncating $m^{1/8}$ singular values
leaves order-one normalized nuclear energy. The full matrix even has
every entry at least $1/2$. Sparsity of the \emph{heavy part}, together
with its actual conditional denominator, is the obstruction.
The table also permits local rank-one upper completions with small
average error: when all four heavy entries vanish, the matrix already
has rank one; otherwise the constant matrix $3J/2$ is an upper
completion, with total error at most four. The expectation of that
error is at most $16/t$. Such independently chosen local completions
need not be the coherent completions of an actual circuit.

\section{Binding the new finite check to an actual cut}
The replay uses $t=4$, not an admissible-large-graph simulation.
There are $k=4$ row bits and $k=4$ column bits, with two two-variable
source groups for each bit. Thus there are sixteen groups and
thirty-two variables. Both options are designated nonroot variables.
For a row string $u$, let $r_u$ choose that string in both row groups
per bit. For a column string $v$, let $c_v$ choose it only in the
first column group per bit and $d_v$ choose it in both column groups.
Let $S$ contain all eight row groups, and define
\[
 V=\sum_v d_v=\prod_{i=1}^k(z_i z_i'+w_i w_i'),\quad
 P=\sum_{u,v:B_{uv}=1}r_u c_v,\quad
 A=\tfrac12F_S+P,\quad Q=AV.                            \tag{RG13}
\]
Build $V$ and $P$ before $A$, then the final output product.
$G(A)=S$, $G(V)=\overline S$, and $A$ is the last balanced gate.
The compulsory cut therefore has divided context exactly $V$;
cutting $A$ kills the output. The replay verifies this using the
existing exact shared-DAG extractor, not an asserted product split.

The full Boolean cube satisfies $Q\le3F/2$. Indeed,
$P\le(\sum_u r_u)(\sum_v c_v)$, $\sum_u r_u\le F_S$,
and
\[
 \left(\sum_v c_v\right)V\le F_{\overline S},
 \qquad V\le F_{\overline S}.
\]
For the first inequality, at each column bit use
$z_i z_i'+w_i w_i'\le z_i'+w_i'$ on Boolean inputs, and multiply
by $z_i+w_i$. Zero faces cause no exception.
For every cutoff $0\le h<k$, precisely $PV$ is heavy: it has
$k$ optional nonroot groups in the first factor, while
$F_SV/2$ has none. On the correlated bit corners its matrices
are exactly (RG9), with $a_{\rm heavy}=1/4$, not the full mean $3/4$.
The finite parameters $k_0=3$, $h=1$ meet the strong-good/thick
inequalities; they are not the asymptotic $N$ parameter recipe.

For each fixed row string, a column Boolean box is described by
its available choices $C_i,D_i$ in the two groups.
Put
\[
 v_{\rm box}=\prod_i|C_i\cap D_i|,\quad
 f_{\rm box}=\prod_i|C_i||D_i|,\quad
 n_u(C)=\#\{v:B_{uv}=1,\ v_i\in C_i\}.
\]
Its defect is exactly
$n_u(C)v_{\rm box}-n_u(C\cap D)$.
An arbitrary row box is a positive sum of these row patterns,
of total weight at most $F_S$. Thus enumerating the $3^8$
nonempty column boxes and sixteen consistent rows, together with
formal coverage of zero faces, determines the full-cube maximum
$b_Q$ exactly. This is a proved reduction, not individual enumeration
of all $2^{32}$ Boolean inputs.

For a sparse flip outcome, let ${\cal U},{\cal V}$ be the row and
column bit subcubes spanned coordinatewise by the two strings.
The actual local K1 completion at a row $u$ is
$1/2+n_u({\cal V})$, independent of its column. Moreover
\[
 D_Q(U_j)=(|{\cal V}|-1)
           \sum_{u\in{\cal U},v\in{\cal V}}B_{uv}.       \tag{RG14}
\]
The checker verifies every one of the $16^4=65536$ base/flip outcomes,
including repeated corners, against these actual completions.
It retains the original conditional bit-flip weights $p=1/8$.
The exact mean raw defect is
\[
 \E D_Q(U_j)=\frac1t(1+p)^k
                    \bigl((1+3p)^k-(1+p)^k\bigr),      \tag{RG15}
\]
and the union normalization is $(1+3p)^{2k}$.
Equation (RG15) follows by first fixing the masks: each of the
$2^{|e_r|+|e_c|}$ subcube corners has incidence probability $1/t$.
In the abstract family with $q\ge k$ and $p=1/(2q)$ this is at most
$e^2/t$. Even the actual \emph{conditional} local defect budget can
therefore vanish while (RG10) stays bounded away from zero.

This observation is not a full-sampler counterexample. These circuits
have an explicit $O(kt^3)$ representation, not a proved polynomial
size in $k$, and the matrix concerns a designated conditional history.
Also $b_Q\ge1/2$: select an edge of the column hypercube across the
boundary of a nonempty proper incidence row, fix the second groups to
the nonincident endpoint, and allow both choices in just that first
group. The defect is one and $F=2$. Thus neither a small whole-output
defect nor the original outer average required to refute (RG8)
has been obtained. A polynomial-size exponential remainder cannot
be dropped on the basis of this diagnostic.

For the finite case the full Gram spectrum is
$144^1,4^{12},0^3$; hence $r=1$, $\tau=24$ and
$\tau^2/(m^3a_{\rm heavy})=9/16$.
The exact full-cube maximum is $b_Q=3/4$.
The original circuit has 320 arithmetic gates, 5120 formal monomials,
and 4096 projected monomials; only one cut term is nonzero.
The inspected Jacobian operation checks that exact complete spectrum.
The local replay independently checks all shifted nullities and binds
the payload and returned matrix to the actual circuit values.
Run \texttt{python3 research/boolean-local-signed-stability/%
local\_rank\_gluing\_check.py} from the campaign root to replay the new
checks. It reads the dedicated typed payload and output but also
recomputes the matrix ranks independently over rational numbers.
No earlier certificate is rerun, and no Lean compilation is claimed.

\paragraph{Remaining mathematical task.}
Retire only the generic matrix-only implication. The actual-circuit,
full-original-sampler conversion
$\mathcal T(Q)\le C_0\sum_j\E_\mu\phi(Z_j)
+K M_n^\alpha(s+1)^\beta e^{-cN}$ remains a possible sufficient route;
so does a direct proof of (RG8) using more of the defect than minors.
Neither is proved here. Any successful continuation must exploit
circuit structure and the original outer averaging, or introduce a
different genuine spectral modulus. No entrywise mass substitute,
unchanged failed pointwise shortcut, or finite diagnostic completes
the requested fixed-factor theorem.
\end{document}
